% --- Archon physics formalization source begin ---
% archon:physics
% archon:covers PhyXMiniProblems/problem_phyx_mini_0883.lean
% archon:source-report reports/phyx_mini/problem_phyx_mini_0883.source.json

\chapter{Physics problem phyx\_mini\_0883}
\label{ch:PhyXMiniProblems_problem_phyx_mini_0883}

\paragraph{Problem source.}
\#\# Physical scenario

The metal wire in figure moves with speed \textbackslash{}( v \textbackslash{}) parallel to a straight wire that is carrying current \textbackslash{}( I \textbackslash{}). The distance between the two wires is \textbackslash{}( d \textbackslash{}).

\#\# Question

Find an expression for the potential difference between the two ends of the moving wire.

\#\# Answer choices

- A: A: \textbackslash{}( (\textbackslash{}mu v\_0 I /\textbackslash{}pi) \textbackslash{}ln[(d+l)/d] \textbackslash{})

- B: B: \textbackslash{}( (\textbackslash{}mu v\_0 I / 3\textbackslash{}pi) \textbackslash{}ln[(d+l)/d] \textbackslash{})

- C: C: \textbackslash{}( (\textbackslash{}mu v\_0 I / 2\textbackslash{}pi) \textbackslash{}ln[(d+l)/d] \textbackslash{})

- D: D: \textbackslash{}( (\textbackslash{}mu v\_0 I / 6\textbackslash{}pi) \textbackslash{}ln[(d+l)/d] \textbackslash{})

\#\# Figure caption (auxiliary; use the image as primary evidence)

The image depicts a physical setup often used in physics to illustrate the principles of electromagnetic interactions:

1. **Objects:**
   - A vertical orange wire labeled with \textbackslash{}( I \textbackslash{}) on the left side, indicating the direction of the current flowing upward.
   - A horizontal gray rectangular bar on the right side from which an upward-pointing green arrow labeled \textbackslash{}( v \textbackslash{}) is emanating, indicating velocity.
   
2. **Text and Labels:**
   - \textbackslash{}( I \textbackslash{}) is labeled next to the wire, pointing upward, representing current flow.
   - \textbackslash{}( v \textbackslash{}) is labeled next to the green arrow above the rectangular bar, representing velocity.
   - The distance \textbackslash{}( d \textbackslash{}) is marked with a double-headed arrow between the wire and rectangle.
   - The length \textbackslash{}( l \textbackslash{}) is marked with a double-headed arrow along the bar.

3. **Relationships:**
   - The wire carrying the current \textbackslash{}( I \textbackslash{}) is oriented vertically.
   - The bar moves perpendicular to the wire with velocity \textbackslash{}( v \textbackslash{}).
   - The distance \textbackslash{}( d \textbackslash{}) represents the separation between the wire and the beginning of the rectangular bar.
   - The length \textbackslash{}( l \textbackslash{}) indicates the length of the rectangular bar.

This diagram is likely used to discuss concepts like magnetic fields, electromotive force (EMF), or Lorentz force in the context of moving conductors in a magnetic field.

\#\# Dataset metadata

Electromagnetic Induction; Physical Model Grounding Reasoning, Multi-Formula Reasoning

\paragraph{Recorded answer/context.}
C: C: \textbackslash{}( (\textbackslash{}mu v\_0 I / 2\textbackslash{}pi) \textbackslash{}ln[(d+l)/d] \textbackslash{})

\paragraph{Figure/image path.}
/root/proposal\_for\_physic/science-mango/phyx\_mini\_run/phyx\_data/test\_image/883.png

\paragraph{Formalization target.}
create a compiling Lean file with sorry bodies at `PhyXMiniProblems/problem\_phyx\_mini\_0883.lean`. The Lean declarations must preserve the physical quantities, dimensions or dimensional roles, figure labels, governing-law hypotheses, and final relation expressed by this problem.
Use Mathlib/Physlib names found through LeanExplore where available. If a physics API is missing, introduce faithful local abstractions rather than scalar placeholder aliases.

\begin{theorem}[Physics formalization target]
\label{thm:physics:phyx_mini_0883:target}
\lean{PhyXMiniProblems.ProblemPhyXMini0883.nearEndMinusFarEndPotential_eq_answer_C}
\uses{def:physics:phyx-mini-0883:phyxminiproblems-problemphyxmini0883-lengthinmeters, def:physics:phyx-mini-0883:phyxminiproblems-problemphyxmini0883-speedinmeterspersecond, def:physics:phyx-mini-0883:phyxminiproblems-problemphyxmini0883-currentinamperes, def:physics:phyx-mini-0883:phyxminiproblems-problemphyxmini0883-permeabilityinhenriespermeter, def:physics:phyx-mini-0883:phyxminiproblems-problemphyxmini0883-potentialdifferenceinvolts, def:physics:phyx-mini-0883:phyxminiproblems-problemphyxmini0883-straightwiremovingbarsetup, def:physics:phyx-mini-0883:phyxminiproblems-problemphyxmini0883-matchesstraightwiremovingbarscenario, def:physics:phyx-mini-0883:phyxminiproblems-problemphyxmini0883-matchessuppliedstraightwiremovingbarfigure, def:physics:phyx-mini-0883:phyxminiproblems-problemphyxmini0883-satisfiesstraightwiremotionalemflaws}
The assigned autoformalize agent should translate this physics problem into Lean declarations in the covered file, with theorem and lemma proof bodies written as `by sorry`.
\end{theorem}
\begin{proof}
This is an autoformalization task, not a proof task. Produce statements that can later be proved by the physics prover without weakening the physical model.
\end{proof}

% --- Archon declaration topology begin ---
\section{Declaration topology}
\begin{definition}[speed Dimension]
  \label{def:physics:phyx-mini-0883:phyxminiproblems-problemphyxmini0883-speeddimension}
  \lean{PhyXMiniProblems.ProblemPhyXMini0883.speedDimension}
  Speed has dimension L T⁻¹
\end{definition}
\begin{proof}
  This is the declared model definition of speed Dimension.
\end{proof}
\begin{definition}[electric Current Dimension]
  \label{def:physics:phyx-mini-0883:phyxminiproblems-problemphyxmini0883-electriccurrentdimension}
  \lean{PhyXMiniProblems.ProblemPhyXMini0883.electricCurrentDimension}
  Electric current has dimension charge per time, C T⁻¹
\end{definition}
\begin{proof}
  This is the declared model definition of electric Current Dimension.
\end{proof}
\begin{definition}[magnetic Permeability Dimension]
  \label{def:physics:phyx-mini-0883:phyxminiproblems-problemphyxmini0883-magneticpermeabilitydimension}
  \lean{PhyXMiniProblems.ProblemPhyXMini0883.magneticPermeabilityDimension}
  Magnetic permeability has SI dimension M L C⁻²
\end{definition}
\begin{proof}
  This is the declared model definition of magnetic Permeability Dimension.
\end{proof}
\begin{definition}[magnetic Flux Density Dimension]
  \label{def:physics:phyx-mini-0883:phyxminiproblems-problemphyxmini0883-magneticfluxdensitydimension}
  \lean{PhyXMiniProblems.ProblemPhyXMini0883.magneticFluxDensityDimension}
  Magnetic flux density has SI dimension M T⁻¹ C⁻¹
\end{definition}
\begin{proof}
  This is the declared model definition of magnetic Flux Density Dimension.
\end{proof}
\begin{definition}[electric Field Dimension]
  \label{def:physics:phyx-mini-0883:phyxminiproblems-problemphyxmini0883-electricfielddimension}
  \lean{PhyXMiniProblems.ProblemPhyXMini0883.electricFieldDimension}
  Motional-emf density has the electric-field dimension M L T⁻² C⁻¹
\end{definition}
\begin{proof}
  This is the declared model definition of electric Field Dimension.
\end{proof}
\begin{definition}[electric Potential Dimension]
  \label{def:physics:phyx-mini-0883:phyxminiproblems-problemphyxmini0883-electricpotentialdimension}
  \lean{PhyXMiniProblems.ProblemPhyXMini0883.electricPotentialDimension}
  Electric potential difference has dimension M L² T⁻² C⁻¹
\end{definition}
\begin{proof}
  This is the declared model definition of electric Potential Dimension.
\end{proof}
\begin{definition}[Length Quantity]
  \label{def:physics:phyx-mini-0883:phyxminiproblems-problemphyxmini0883-lengthquantity}
  \lean{PhyXMiniProblems.ProblemPhyXMini0883.LengthQuantity}
  A nonnegative, unit-independent physical length
\end{definition}
\begin{proof}
  This is the declared model definition of Length Quantity.
\end{proof}
\begin{definition}[Speed Quantity]
  \label{def:physics:phyx-mini-0883:phyxminiproblems-problemphyxmini0883-speedquantity}
  \lean{PhyXMiniProblems.ProblemPhyXMini0883.SpeedQuantity}
  \uses{def:physics:phyx-mini-0883:phyxminiproblems-problemphyxmini0883-speeddimension}
  A nonnegative, unit-independent physical speed
\end{definition}
\begin{proof}
  This is the declared model definition of Speed Quantity.
\end{proof}
\begin{definition}[Electric Current Magnitude]
  \label{def:physics:phyx-mini-0883:phyxminiproblems-problemphyxmini0883-electriccurrentmagnitude}
  \lean{PhyXMiniProblems.ProblemPhyXMini0883.ElectricCurrentMagnitude}
  \uses{def:physics:phyx-mini-0883:phyxminiproblems-problemphyxmini0883-electriccurrentdimension}
  A nonnegative electric-current magnitude
\end{definition}
\begin{proof}
  This is the declared model definition of Electric Current Magnitude.
\end{proof}
\begin{definition}[Magnetic Permeability Quantity]
  \label{def:physics:phyx-mini-0883:phyxminiproblems-problemphyxmini0883-magneticpermeabilityquantity}
  \lean{PhyXMiniProblems.ProblemPhyXMini0883.MagneticPermeabilityQuantity}
  \uses{def:physics:phyx-mini-0883:phyxminiproblems-problemphyxmini0883-magneticpermeabilitydimension}
  A nonnegative magnetic permeability
\end{definition}
\begin{proof}
  This is the declared model definition of Magnetic Permeability Quantity.
\end{proof}
\begin{definition}[Magnetic Flux Density Magnitude]
  \label{def:physics:phyx-mini-0883:phyxminiproblems-problemphyxmini0883-magneticfluxdensitymagnitude}
  \lean{PhyXMiniProblems.ProblemPhyXMini0883.MagneticFluxDensityMagnitude}
  \uses{def:physics:phyx-mini-0883:phyxminiproblems-problemphyxmini0883-magneticfluxdensitydimension}
  A nonnegative magnetic-flux-density magnitude
\end{definition}
\begin{proof}
  This is the declared model definition of Magnetic Flux Density Magnitude.
\end{proof}
\begin{definition}[Motional Emf Density Magnitude]
  \label{def:physics:phyx-mini-0883:phyxminiproblems-problemphyxmini0883-motionalemfdensitymagnitude}
  \lean{PhyXMiniProblems.ProblemPhyXMini0883.MotionalEmfDensityMagnitude}
  \uses{def:physics:phyx-mini-0883:phyxminiproblems-problemphyxmini0883-electricfielddimension}
  A nonnegative local motional-emf density, measured in volts per metre
\end{definition}
\begin{proof}
  This is the declared model definition of Motional Emf Density Magnitude.
\end{proof}
\begin{definition}[Electric Potential Difference]
  \label{def:physics:phyx-mini-0883:phyxminiproblems-problemphyxmini0883-electricpotentialdifference}
  \lean{PhyXMiniProblems.ProblemPhyXMini0883.ElectricPotentialDifference}
  \uses{def:physics:phyx-mini-0883:phyxminiproblems-problemphyxmini0883-electricpotentialdimension}
  A nonnegative endpoint potential difference
\end{definition}
\begin{proof}
  This is the declared model definition of Electric Potential Difference.
\end{proof}
\begin{definition}[nonnegative SIReadout]
  \label{def:physics:phyx-mini-0883:phyxminiproblems-problemphyxmini0883-nonnegativesireadout}
  \lean{PhyXMiniProblems.ProblemPhyXMini0883.nonnegativeSIReadout}
  Read a nonnegative dimensionful quantity in coherent SI units
\end{definition}
\begin{proof}
  This is the declared model definition of nonnegative SIReadout.
\end{proof}
\begin{definition}[length In Meters]
  \label{def:physics:phyx-mini-0883:phyxminiproblems-problemphyxmini0883-lengthinmeters}
  \lean{PhyXMiniProblems.ProblemPhyXMini0883.lengthInMeters}
  \uses{def:physics:phyx-mini-0883:phyxminiproblems-problemphyxmini0883-lengthquantity, def:physics:phyx-mini-0883:phyxminiproblems-problemphyxmini0883-nonnegativesireadout}
  Read a length in metres
\end{definition}
\begin{proof}
  This is the declared model definition of length In Meters.
\end{proof}
\begin{definition}[speed In Meters Per Second]
  \label{def:physics:phyx-mini-0883:phyxminiproblems-problemphyxmini0883-speedinmeterspersecond}
  \lean{PhyXMiniProblems.ProblemPhyXMini0883.speedInMetersPerSecond}
  \uses{def:physics:phyx-mini-0883:phyxminiproblems-problemphyxmini0883-speedquantity, def:physics:phyx-mini-0883:phyxminiproblems-problemphyxmini0883-nonnegativesireadout}
  Read a speed in metres per second
\end{definition}
\begin{proof}
  This is the declared model definition of speed In Meters Per Second.
\end{proof}
\begin{definition}[current In Amperes]
  \label{def:physics:phyx-mini-0883:phyxminiproblems-problemphyxmini0883-currentinamperes}
  \lean{PhyXMiniProblems.ProblemPhyXMini0883.currentInAmperes}
  \uses{def:physics:phyx-mini-0883:phyxminiproblems-problemphyxmini0883-electriccurrentmagnitude, def:physics:phyx-mini-0883:phyxminiproblems-problemphyxmini0883-nonnegativesireadout}
  Read an electric-current magnitude in amperes
\end{definition}
\begin{proof}
  This is the declared model definition of current In Amperes.
\end{proof}
\begin{definition}[permeability In Henries Per Meter]
  \label{def:physics:phyx-mini-0883:phyxminiproblems-problemphyxmini0883-permeabilityinhenriespermeter}
  \lean{PhyXMiniProblems.ProblemPhyXMini0883.permeabilityInHenriesPerMeter}
  \uses{def:physics:phyx-mini-0883:phyxminiproblems-problemphyxmini0883-magneticpermeabilityquantity, def:physics:phyx-mini-0883:phyxminiproblems-problemphyxmini0883-nonnegativesireadout}
  Read magnetic permeability in henries per metre
\end{definition}
\begin{proof}
  This is the declared model definition of permeability In Henries Per Meter.
\end{proof}
\begin{definition}[magnetic Flux Density In Teslas]
  \label{def:physics:phyx-mini-0883:phyxminiproblems-problemphyxmini0883-magneticfluxdensityinteslas}
  \lean{PhyXMiniProblems.ProblemPhyXMini0883.magneticFluxDensityInTeslas}
  \uses{def:physics:phyx-mini-0883:phyxminiproblems-problemphyxmini0883-magneticfluxdensitymagnitude, def:physics:phyx-mini-0883:phyxminiproblems-problemphyxmini0883-nonnegativesireadout}
  Read a magnetic-flux-density magnitude in teslas
\end{definition}
\begin{proof}
  This is the declared model definition of magnetic Flux Density In Teslas.
\end{proof}
\begin{definition}[motional Emf Density In Volts Per Meter]
  \label{def:physics:phyx-mini-0883:phyxminiproblems-problemphyxmini0883-motionalemfdensityinvoltspermeter}
  \lean{PhyXMiniProblems.ProblemPhyXMini0883.motionalEmfDensityInVoltsPerMeter}
  \uses{def:physics:phyx-mini-0883:phyxminiproblems-problemphyxmini0883-motionalemfdensitymagnitude, def:physics:phyx-mini-0883:phyxminiproblems-problemphyxmini0883-nonnegativesireadout}
  Read local motional-emf density in volts per metre
\end{definition}
\begin{proof}
  This is the declared model definition of motional Emf Density In Volts Per Meter.
\end{proof}
\begin{definition}[potential Difference In Volts]
  \label{def:physics:phyx-mini-0883:phyxminiproblems-problemphyxmini0883-potentialdifferenceinvolts}
  \lean{PhyXMiniProblems.ProblemPhyXMini0883.potentialDifferenceInVolts}
  \uses{def:physics:phyx-mini-0883:phyxminiproblems-problemphyxmini0883-electricpotentialdifference, def:physics:phyx-mini-0883:phyxminiproblems-problemphyxmini0883-nonnegativesireadout}
  Read an electric potential difference in volts
\end{definition}
\begin{proof}
  This is the declared model definition of potential Difference In Volts.
\end{proof}
\begin{definition}[Source Wire Model]
  \label{def:physics:phyx-mini-0883:phyxminiproblems-problemphyxmini0883-sourcewiremodel}
  \lean{PhyXMiniProblems.ProblemPhyXMini0883.SourceWireModel}
  Idealization assigned to the current-carrying source wire
\end{definition}
\begin{proof}
  This is the declared model definition of Source Wire Model.
\end{proof}
\begin{definition}[Moving Bar Model]
  \label{def:physics:phyx-mini-0883:phyxminiproblems-problemphyxmini0883-movingbarmodel}
  \lean{PhyXMiniProblems.ProblemPhyXMini0883.MovingBarModel}
  Idealization assigned to the moving metal object
\end{definition}
\begin{proof}
  This is the declared model definition of Moving Bar Model.
\end{proof}
\begin{definition}[Bar Endpoint]
  \label{def:physics:phyx-mini-0883:phyxminiproblems-problemphyxmini0883-barendpoint}
  \lean{PhyXMiniProblems.ProblemPhyXMini0883.BarEndpoint}
  The two ends of the moving bar, distinguished by distance from the wire
\end{definition}
\begin{proof}
  This is the declared model definition of Bar Endpoint.
\end{proof}
\begin{definition}[Figure Quantity Label]
  \label{def:physics:phyx-mini-0883:phyxminiproblems-problemphyxmini0883-figurequantitylabel}
  \lean{PhyXMiniProblems.ProblemPhyXMini0883.FigureQuantityLabel}
  The four physical quantities explicitly labelled in the raster
\end{definition}
\begin{proof}
  This is the declared model definition of Figure Quantity Label.
\end{proof}
\begin{definition}[expected Printed Symbol]
  \label{def:physics:phyx-mini-0883:phyxminiproblems-problemphyxmini0883-expectedprintedsymbol}
  \lean{PhyXMiniProblems.ProblemPhyXMini0883.expectedPrintedSymbol}
  \uses{def:physics:phyx-mini-0883:phyxminiproblems-problemphyxmini0883-figurequantitylabel}
  The printed symbol expected for each labelled quantity
\end{definition}
\begin{proof}
  This is the declared model definition of expected Printed Symbol.
\end{proof}
\begin{definition}[Figure Axis]
  \label{def:physics:phyx-mini-0883:phyxminiproblems-problemphyxmini0883-figureaxis}
  \lean{PhyXMiniProblems.ProblemPhyXMini0883.FigureAxis}
  The two axes needed to describe the idealized planar figure
\end{definition}
\begin{proof}
  This is the declared model definition of Figure Axis.
\end{proof}
\begin{definition}[Figure Arrow Direction]
  \label{def:physics:phyx-mini-0883:phyxminiproblems-problemphyxmini0883-figurearrowdirection}
  \lean{PhyXMiniProblems.ProblemPhyXMini0883.FigureArrowDirection}
  Directions of the two single-headed arrows visible in the figure
\end{definition}
\begin{proof}
  This is the declared model definition of Figure Arrow Direction.
\end{proof}
\begin{definition}[Straight Wire Moving Bar Figure]
  \label{def:physics:phyx-mini-0883:phyxminiproblems-problemphyxmini0883-straightwiremovingbarfigure}
  \lean{PhyXMiniProblems.ProblemPhyXMini0883.StraightWireMovingBarFigure}
  \uses{def:physics:phyx-mini-0883:phyxminiproblems-problemphyxmini0883-figurequantitylabel, def:physics:phyx-mini-0883:phyxminiproblems-problemphyxmini0883-figureaxis, def:physics:phyx-mini-0883:phyxminiproblems-problemphyxmini0883-figurearrowdirection}
  Literal and idealized readouts of image 883.png. Quantitative distances are kept in the physical setup; this structure records only presentation and orientation facts from the raster
\end{definition}
\begin{proof}
  This is the declared model definition of Straight Wire Moving Bar Figure.
\end{proof}
\begin{definition}[Straight Wire Moving Bar Setup]
  \label{def:physics:phyx-mini-0883:phyxminiproblems-problemphyxmini0883-straightwiremovingbarsetup}
  \lean{PhyXMiniProblems.ProblemPhyXMini0883.StraightWireMovingBarSetup}
  \uses{def:physics:phyx-mini-0883:phyxminiproblems-problemphyxmini0883-lengthquantity, def:physics:phyx-mini-0883:phyxminiproblems-problemphyxmini0883-speedquantity, def:physics:phyx-mini-0883:phyxminiproblems-problemphyxmini0883-electriccurrentmagnitude, def:physics:phyx-mini-0883:phyxminiproblems-problemphyxmini0883-magneticpermeabilityquantity, def:physics:phyx-mini-0883:phyxminiproblems-problemphyxmini0883-magneticfluxdensitymagnitude, def:physics:phyx-mini-0883:phyxminiproblems-problemphyxmini0883-motionalemfdensitymagnitude, def:physics:phyx-mini-0883:phyxminiproblems-problemphyxmini0883-electricpotentialdifference, def:physics:phyx-mini-0883:phyxminiproblems-problemphyxmini0883-sourcewiremodel, def:physics:phyx-mini-0883:phyxminiproblems-problemphyxmini0883-movingbarmodel, def:physics:phyx-mini-0883:phyxminiproblems-problemphyxmini0883-barendpoint, def:physics:phyx-mini-0883:phyxminiproblems-problemphyxmini0883-straightwiremovingbarfigure}
  Independent physical quantities in the apparatus. The field profiles and the potential difference are not defined from the answer; the governing laws that relate them are separate assumptions below
\end{definition}
\begin{proof}
  This is the declared model definition of Straight Wire Moving Bar Setup.
\end{proof}
\begin{definition}[Matches Straight Wire Moving Bar Scenario]
  \label{def:physics:phyx-mini-0883:phyxminiproblems-problemphyxmini0883-matchesstraightwiremovingbarscenario}
  \lean{PhyXMiniProblems.ProblemPhyXMini0883.MatchesStraightWireMovingBarScenario}
  \uses{def:physics:phyx-mini-0883:phyxminiproblems-problemphyxmini0883-lengthinmeters, def:physics:phyx-mini-0883:phyxminiproblems-problemphyxmini0883-speedinmeterspersecond, def:physics:phyx-mini-0883:phyxminiproblems-problemphyxmini0883-currentinamperes, def:physics:phyx-mini-0883:phyxminiproblems-problemphyxmini0883-permeabilityinhenriespermeter, def:physics:phyx-mini-0883:phyxminiproblems-problemphyxmini0883-straightwiremovingbarsetup}
  The component idealizations and nondegenerate physical magnitudes
\end{definition}
\begin{proof}
  This is the declared model definition of Matches Straight Wire Moving Bar Scenario.
\end{proof}
\begin{definition}[Matches Supplied Straight Wire Moving Bar Figure]
  \label{def:physics:phyx-mini-0883:phyxminiproblems-problemphyxmini0883-matchessuppliedstraightwiremovingbarfigure}
  \lean{PhyXMiniProblems.ProblemPhyXMini0883.MatchesSuppliedStraightWireMovingBarFigure}
  \uses{def:physics:phyx-mini-0883:phyxminiproblems-problemphyxmini0883-lengthinmeters, def:physics:phyx-mini-0883:phyxminiproblems-problemphyxmini0883-expectedprintedsymbol, def:physics:phyx-mini-0883:phyxminiproblems-problemphyxmini0883-straightwiremovingbarsetup}
  The labels, orientations, arrow directions, and distance assignments supplied by the primary raster. In particular, the current and velocity arrows are both upward (parallel), while the bar is horizontal. The far-end radius is d + l; no potential-difference value is asserted here
\end{definition}
\begin{proof}
  This is the declared model definition of Matches Supplied Straight Wire Moving Bar Figure.
\end{proof}
\begin{definition}[Satisfies Straight Wire Motional Emf Laws]
  \label{def:physics:phyx-mini-0883:phyxminiproblems-problemphyxmini0883-satisfiesstraightwiremotionalemflaws}
  \lean{PhyXMiniProblems.ProblemPhyXMini0883.SatisfiesStraightWireMotionalEmfLaws}
  \uses{def:physics:phyx-mini-0883:phyxminiproblems-problemphyxmini0883-lengthinmeters, def:physics:phyx-mini-0883:phyxminiproblems-problemphyxmini0883-speedinmeterspersecond, def:physics:phyx-mini-0883:phyxminiproblems-problemphyxmini0883-currentinamperes, def:physics:phyx-mini-0883:phyxminiproblems-problemphyxmini0883-permeabilityinhenriespermeter, def:physics:phyx-mini-0883:phyxminiproblems-problemphyxmini0883-magneticfluxdensityinteslas, def:physics:phyx-mini-0883:phyxminiproblems-problemphyxmini0883-motionalemfdensityinvoltspermeter, def:physics:phyx-mini-0883:phyxminiproblems-problemphyxmini0883-potentialdifferenceinvolts, def:physics:phyx-mini-0883:phyxminiproblems-problemphyxmini0883-straightwiremovingbarsetup}
  The governing electromagnetic laws used in the intended derivation. The first field is the magnitude B(r) = μ I / (2 π r) around an ideal infinite straight wire. For the upward velocity and current shown in the figure, the magnitude of the radially oriented motional-emf density along the bar is v B(r). The near-end-minus-far-end potential is the line integral of that density from the near radius to the far radius. These are general local and integral laws.
\end{definition}
\begin{proof}
  This is the declared model definition of Satisfies Straight Wire Motional Emf Laws.
\end{proof}
% --- Archon declaration topology end ---
% --- Archon physics formalization source end ---
